% ID: SEL-D0005-OB-E02
% ESTADO: validada. S006, 2026-11-07; tres pasadas del mismo autor y revisión matemática, didáctica y visual.
\textbf{a) Sustitución del precio.}\par\nopagebreak[4]
El precio 0.50 euros/kg pertenece al dominio $[0,4]$. Sustituyendo,
\[
D(0.5)=-200(0.5)^3+2100(0.5)^2-7200(0.5)+10000
=\boxed{6900\text{ kg}}.
\]

\textbf{b) Mínimo absoluto en el dominio.}\par\nopagebreak[4]
La función es continua en $[0,4]$, por lo que alcanza mínimo. Debemos comparar los puntos estacionarios interiores y los extremos del intervalo:
\[
D'(x)=-600x^2+4200x-7200=-600(x-3)(x-4).
\]
Los ceros son 3 y4, pero 4 es extremo del dominio, no punto interior. La derivada es negativa en $(0,3)$ y positiva en $(3,4)$: la demanda decrece hasta 3 y después crece hasta 4. Comparamos valores:
\[
D(0)=10000,\qquad D(3)=1900,\qquad D(4)=2000.
\]
Por tanto, \boxed{x=3\text{ euros/kg}} produce una demanda mínima de \boxed{1900\text{ kg}}. No se pide minimizar el precio ni maximizar ingresos $xD(x)$.

Comprobación global alternativa por factorización:
\[
D(x)-D(3)=200(x-3)^2(4.5-x)\geq0\qquad(0\leq x\leq4).
\]
El último factor es positivo en el dominio, y la igualdad exige $x=3$, por lo que el mínimo es único. El modelo aumenta ligeramente entre 3 y 4: se conserva tal como está dado, sin imponer una monotonía económica ajena a sus datos.
